\section{从 RTL 到可实现约束}

RTL 仿真回答“在给定激励下，描述的状态怎样变化”，但芯片或 FPGA 还必须回答“这些变化能否在真实时钟和引脚边界内发生”。工具首先进行展开与静态检查：解析参数和 generate 分支，确定端口宽度与驱动关系，再报告未连接端口、位宽截断、组合环、多重驱动、未完全赋值以及不可达状态。lint 不是语法检查的别名；它要在仿真尚未表现异常时指出可能形成锁存器、X 扩散或硬件歧义的结构。

时序约束把外部契约写给静态时序分析器。主时钟约束给出周期与波形；派生时钟说明分频、倍频或门控后的相位关系；输入延迟描述外部发送器相对本芯片采样沿何时给出数据；输出延迟描述板级接收器要求本芯片何时稳定输出。没有输入/输出延迟的设计即使内部路径全部通过，也没有证明芯片边界满足接口时序。

\begin{longtable}{L{0.18\linewidth}L{0.34\linewidth}L{0.38\linewidth}}
\toprule
约束对象 & 它表达的硬件事实 & 缺失或误用的后果 \\
\midrule
主时钟 & 寄存器采样边沿和周期预算 & 工具不知道组合路径必须多快到达 \\
派生时钟 & 两个相关时钟的频率、相位和来源 & 把相关路径误当异步，或使用错误采样沿 \\
输入延迟 & 板外器件已经消耗的周期预算 & 内部输入路径通过，系统级建立时间仍可能失败 \\
输出延迟 & 板外接收器要求保留的建立/保持预算 & 输出到达过晚或变化过早 \\
时钟不确定度 & 抖动、偏斜和建模余量 & 纸面周期没有给真实分配网络留出裕量 \\
例外路径 & 功能上不按普通单周期采样的路径 & 错误例外会隐藏真实失败，缺少例外会产生伪告警 \\
\bottomrule
\end{longtable}

false path 只能用于功能上永远不会被同步采样的路径，例如已经由异步协议隔离的跨域关系；multicycle path 表示接收端确实允许多个周期后采样，因此还要同步修正保持检查。两者都不是“让报告变绿”的选项。每条例外都应能回到电路模式、状态条件和接口协议，证明被排除的采样事件不会发生。

\section{静态时序分析把路径拆成预算}

静态时序分析（Static Timing Analysis, STA）不依赖某一组输入向量，而是沿时序图检查所有可达寄存器路径。建立时间检查比较数据最晚到达时刻与接收寄存器的要求时刻；保持时间检查比较数据最早到达时刻与旧值必须维持到的时刻。两者使用不同延迟角：建立失败通常来自数据路径过慢，保持失败通常来自数据路径过快或时钟偏斜不利。

设发射寄存器时钟到输出为 $t_{cq}$，组合路径最大延迟为 $t_{comb,max}$，接收寄存器建立时间为 $t_{setup}$，时钟偏斜与不确定度预算合计为 $t_u$，则单周期路径至少要求

\[
T_{clk} \ge t_{cq}+t_{comb,max}+t_{setup}+t_u.
\]

若 $T_{clk}=2.0$ ns，四项分别为 $0.12$、$1.62$、$0.08$ 和 $0.10$ ns，要求时刻与到达时刻之差只有 $0.08$ ns。这个正裕量不是“还有很多空间”，因为布线拥塞、工艺角和串扰模型更新都可能消耗它。报告阅读应先确认起点、终点、时钟关系和例外是否正确，再看组合单元、扇出、布线和逻辑级数；直接优化报告中最长的一串单元，可能只是在修一条约束错误的伪路径。

建立失败可以通过减少逻辑级数、复制高扇出驱动、调整布局或增加流水级修复。保持失败不能靠降低时钟频率解决，因为同一采样沿附近的数据仍然到得过早；常见修复是在数据路径增加受控延迟、调整时钟树或修正不合理的跨域约束。任何自动修复后都要重新检查所有角和所有模式，避免一个角的改动在另一个角制造新失败。

\section{CDC、复位与形式性质}

跨时钟域（Clock Domain Crossing, CDC）检查关注的不是组合功能，而是数据怎样跨越无固定相位关系的采样边界。单 bit 电平可以经两级同步器降低亚稳态传播概率；脉冲必须保证目标域能够看到，或改成 toggle/握手协议；多 bit 数据要用稳定快照握手、异步 FIFO 或带编码保证的专用结构。把多位总线的每一位分别接两级同步器，仍可能在目标域拼出源域从未存在过的混合值。

复位域跨越（Reset Domain Crossing, RDC）还要检查不同模块解除复位的先后关系。异步置位可以快速把输出拉到安全状态，但解除通常要在本地时钟域同步完成。若生产者已经开始发送而消费者仍处于复位，valid、credit 或指针状态会失配。设计应规定复位后谁先声明就绪，以及旧事务如何作废。

形式验证用“所有满足前提的状态和输入”检查性质，而不是只运行若干测试向量。适合 RTL 的性质通常分三类：安全性说明坏事永远不发生，例如 FIFO 不在空时读取；活性说明在公平前提下请求最终得到响应；等价性说明综合或优化前后的可观察行为一致。断言的前提必须描述环境契约，否则工具可能用现实中不可能的输入找到反例，也可能因为前提过强而把错误排除在搜索空间之外。

\begin{longtable}{L{0.18\linewidth}L{0.34\linewidth}L{0.38\linewidth}}
\toprule
验证层 & 主要输入 & 通过时真正证明了什么 \\
\midrule
lint/展开 & RTL 结构、参数和连接 & 未发现选定规则覆盖的静态结构缺陷 \\
仿真 & 测试激励、参考模型、覆盖点 & 给定场景下行为符合预期，未覆盖状态仍未知 \\
形式性质 & 假设、断言和可达状态空间 & 在假设成立的全部可达状态内性质成立 \\
CDC/RDC & 时钟域、复位域和跨域结构 & 跨边界协议符合已识别的安全模板与时序要求 \\
STA & 网表、延迟、时钟和接口约束 & 所有被分析路径在指定模式与工艺角满足时序 \\
实现后检查 & 布局布线、功耗和物理规则 & 逻辑已映射到可制造或可配置资源并满足物理边界 \\
\bottomrule
\end{longtable}

\section{用证据闭合一次 RTL 修改}

一次看似局部的控制器修改可能同时改变状态可达性、组合路径、复位行为和接口背压。完整闭环应从需求中提取不变量与完成条件，先更新断言和测试，再运行 lint、仿真、CDC/RDC、综合与 STA。若修改增加一级流水寄存器，还要同步更新 valid、错误位、事务 ID 和输出延迟；只让数据多走一拍而控制仍按旧拍到达，会形成比原始时序失败更隐蔽的功能错误。

验收记录应保留工具版本、顶层参数、约束集合、未豁免告警、覆盖结果、最差建立/保持裕量以及关键资源变化。这样下一次频率下降、面积上升或跨域告警出现时，才能判断它来自功能改动、约束改动还是实现随机性。第 3 章建立的采样边界在这里成为可检查的时序路径；本章形成的网表与时序证据，则交给后续处理器章节作为“电路确实能按所述拍序运行”的实现基础。

\section*{章末练习}

\begin{longtable}{L{0.18\linewidth}L{0.42\linewidth}L{0.30\linewidth}}
\toprule
任务 & 要完成的产物 & 检查要点 \\
\midrule
约束边界 & 解释为什么只有主时钟约束仍不能证明芯片接口时序 & 输入/输出延迟属于板级契约 \\
例外审查 & 比较 false path 与 multicycle path 的功能含义 & 例外必须由真实采样关系证明 \\
建立计算 & 用正文给出的四项延迟计算 2.0 ns 周期下的裕量 & 到达与要求时刻使用同一单位 \\
保持修复 & 说明为什么降低频率不能修复普通保持失败 & 保持检查围绕同一采样沿 \\
跨域选择 & 为单 bit 电平、窄脉冲和连续多 bit 数据分别选择跨域结构 & 不逐位同步多 bit payload \\
验证闭环 & 为控制器增加一级流水后列出必须同步更新和重新检查的对象 & 数据、控制、错误与事务身份同行 \\
\bottomrule
\end{longtable}

\section*{参考解答与要点}

\begin{longtable}{L{0.18\linewidth}L{0.46\linewidth}L{0.26\linewidth}}
\toprule
任务 & 参考解答 & 必须保留的检查点 \\
\midrule
约束边界 & 主时钟只给出内部采样周期；外部发送器的到达时间和接收器的要求时间还要由输入、输出延迟表达。 & 内部通过不等于板级接口通过 \\
例外审查 & false path 表示功能上不发生同步采样，multicycle 表示允许若干周期后采样且需配套保持关系。 & 不能用例外隐藏普通慢路径 \\
建立计算 & $2.0-(0.12+1.62+0.08+0.10)=0.08$ ns。裕量为正但很小，仍要检查角、模式和建模变化。 & 不遗漏不确定度预算 \\
保持修复 & 保持约束要求新数据不要在同一采样沿后过早到达，改变下一采样沿的周期并不改变这一局部关系。 & 用数据延迟或时钟树修复 \\
跨域选择 & 电平用同步器；窄脉冲用展宽、toggle 或握手；连续多 bit 流用异步 FIFO。 & 协议要定义接受与复位边界 \\
验证闭环 & 更新流水 valid、异常/错误位、ID、接口延迟与断言，重跑仿真、CDC、综合和 STA，并比较覆盖、资源与裕量。 & 延迟变化必须在接口契约中可见 \\
\bottomrule
\end{longtable}
